% !TeX root = part1.tex

\documentclass[../final.tex]{subfiles}

\begin{document}

% ============================================================
% Задача 1
% ============================================================

\newpage
\task[1]{Распад \texorpdfstring{$\varphi$}{φ}-мезона на каоны}

\condition

$\varphi$-мезон движется в лабораторной системе отсчёта
со скоростью $\beta c$ и на лету распадается на пару заряженных
каонов
%
\begin{equation*}
    \varphi \rightarrow K^+ + K^-.
\end{equation*}
%
В системе покоя $\varphi$-мезона оба каона имеют одинаковые
по модулю скорости
%
\begin{equation*}
    v'=\beta'c.
\end{equation*}

Требуется определить максимальный угол, который может составлять
направление движения одного из каонов с направлением движения
$\varphi$-мезона в лабораторной системе, а также максимальный угол
между двумя каонами в лабораторной системе.


% ------------------------------------------------------------
% Схема
% ------------------------------------------------------------

\begin{rbox}[left=18pt,right=18pt,top=14pt,bottom=10pt]{[]}
    \centering

    \begin{tikzpicture}[
        >=Latex,
        every node/.style={text=rtext,font=\small}
    ]

        % =====================================================
        % Система покоя phi-мезона
        % =====================================================

        \begin{scope}[xshift=-3.7cm]

            \coordinate (O1) at (0,0);

            % Оси
            \draw[
                ->,
                draw=rtext,
                line width=0.8pt
            ]
            (-0.4,0) -- (3.1,0)
            node[right] {$x'$};

            \draw[
                ->,
                draw=rtext,
                line width=0.8pt
            ]
            (0,-1.9) -- (0,2.2)
            node[above] {$y'$};

            % Вершина распада
            \fill[rc3] (O1) circle (2.3pt);

            % K+
            \draw[
                ->,
                draw=rtext,
                line width=1pt
            ]
            (O1) -- ++(2.1,1.45);

            \node[above right] at (1.85,1.28)
            {$K^+,\ \beta'c$};

            % K-
            \draw[
                ->,
                draw=rtext,
                line width=1pt
            ]
            (O1) -- ++(-2.1,-1.45);

            \node[below left] at (-1.75,-1.2)
            {$K^-,\ \beta'c$};

            % Угол theta'
            \draw[
                draw=rtext,
                line width=0.8pt
            ]
            (0.85,0)
            arc[
                start angle=0,
                end angle=35,
                radius=0.85
            ];

            \node at (1.08,0.35) {$\theta'$};

            \node[above] at (0,2.45)
            {\textbf{С.О. $\varphi$-мезона}};

        \end{scope}


        % =====================================================
        % Лабораторная система
        % =====================================================

        \begin{scope}[xshift=3.7cm]

            \coordinate (O2) at (0,0);

            % Ось x
            \draw[
                ->,
                draw=rtext,
                line width=0.8pt
            ]
            (-2.8,0) -- (3.2,0)
            node[right] {$x$};

            % Ось y
            \draw[
                ->,
                draw=rtext,
                line width=0.8pt
            ]
            (0,-1.9) -- (0,2.2)
            node[above] {$y$};

            % phi до распада
            \draw[
                ->,
                draw=rc3,
                line width=1.1pt
            ]
            (-2.6,0) -- (O2);

            \node[above] at (-1.35,0)
            {$\varphi,\ \beta c$};

            \fill[rc3] (O2) circle (2.3pt);

            % K+
            \draw[
                ->,
                draw=rtext,
                line width=1pt
            ]
            (O2) -- ++(2.25,1.55);

            \node[above right] at (2.05,1.4) {$K^+$};

            % K-
            \draw[
                ->,
                draw=rtext,
                line width=1pt
            ]
            (O2) -- ++(2.25,-1.25);

            \node[below right] at (2.0,-1.1) {$K^-$};

            % theta_1
            \draw[
                draw=rtext,
                line width=0.8pt
            ]
            (0.85,0)
            arc[
                start angle=0,
                end angle=35,
                radius=0.85
            ];

            \node at (1.08,0.36) {$\theta_1$};

            % theta_2
            \draw[
                draw=rtext,
                line width=0.8pt
            ]
            (0.9,0)
            arc[
                start angle=0,
                end angle=-29,
                radius=0.9
            ];

            \node at (1.12,-0.31) {$\theta_2$};

            % Угол Psi
            \draw[
                draw=rc3,
                line width=0.9pt
            ]
            (1.72,1.18)
            arc[
                start angle=35,
                end angle=-29,
                radius=2.1
            ];

            \node[text=rc3] at (2.15,0.25) {$\Psi$};

            \node[above] at (0,2.45)
            {\textbf{Лабораторная С.О.}};

        \end{scope}

    \end{tikzpicture}

    \captionsetup{
        font=footnotesize,
        labelfont=bf,
        skip=5pt
    }

    \captionof{figure}{
        Распад $\varphi\rightarrow K^+K^-$
        в системе покоя $\varphi$-мезона
        и в лабораторной системе отсчёта.
    }
    \label{fig:phi_to_kaons}
\end{rbox}


\subsection*{Максимальный угол одного каона относительно направления
движения $\varphi$-мезона.}

Направим ось $x$ вдоль скорости $\varphi$-мезона.
Пусть в системе покоя $\varphi$-мезона каон вылетает под углом
$\theta'$ к этой оси.

Используем формулу преобразования угла, полученную ранее:
%
\begin{equation*}
    \tan\theta
    =
    \frac{
        \sin\theta'
    }{
        \gamma
        \left(
            \dfrac{\beta}{\beta'}
            +
            \cos\theta'
        \right)
    },
    \qquad
    \gamma
    =
    \frac{1}{\sqrt{1-\beta^2}}.
\end{equation*}

Сначала рассмотрим случай
%
\begin{equation*}
    \beta>\beta'.
\end{equation*}
%
Введём обозначение
%
\begin{equation*}
    a
    =
    \frac{\beta}{\beta'},
    \qquad
    a>1.
\end{equation*}

Тогда
%
\begin{equation*}
    \tan\theta
    =
    \frac{1}{\gamma}
    \frac{
        \sin\theta'
    }{
        a+\cos\theta'
    }.
\end{equation*}

Максимум угла $\theta$ соответствует максимуму $\tan\theta$.
Дифференцируем:
%
\begin{align*}
    \frac{d}{d\theta'}
    \left(
        \frac{\sin\theta'}{a+\cos\theta'}
    \right)
    &=
    \frac{
        \cos\theta'(a+\cos\theta')
        +
        \sin^2\theta'
    }{
        (a+\cos\theta')^2
    }
    \\[3pt]
    &=
    \frac{
        a\cos\theta'+1
    }{
        (a+\cos\theta')^2
    }.
\end{align*}

Следовательно, условие экстремума имеет вид
%
\begin{equation*}
    a\cos\theta'+1=0,
\end{equation*}
%
откуда
%
\begin{equation*}
    \cos\theta'_{\mathrm{ext}}
    =
    -\frac{1}{a}
    =
    -\frac{\beta'}{\beta}.
\end{equation*}

При этом
%
\begin{equation*}
    \sin\theta'_{\mathrm{ext}}
    =
    \sqrt{
        1-\frac{\beta'^2}{\beta^2}
    }
    =
    \frac{
        \sqrt{\beta^2-\beta'^2}
    }{\beta}.
\end{equation*}

Подставляя это в формулу для угла, получаем
%
\begin{equation*}
    \tan\theta_{\max}
    =
    \frac{
        \beta'
    }{
        \gamma\sqrt{\beta^2-\beta'^2}
    }.
\end{equation*}

Для перехода к синусу используем
%
\begin{equation*}
    \sin\theta
    =
    \frac{\tan\theta}{
        \sqrt{1+\tan^2\theta}
    }.
\end{equation*}

Тогда
%
\begin{equation*}
    1+\tan^2\theta_{\max}
    =
    1+
    \frac{
        \beta'^2
    }{
        \gamma^2(\beta^2-\beta'^2)
    }.
\end{equation*}

С учётом
%
\begin{equation*}
    \gamma^{-2}=1-\beta^2
\end{equation*}
%
после упрощения получаем
%
\begin{equation*}
    \sin\theta_{\max}
    =
    \frac{
        \gamma'\beta'
    }{
        \gamma\beta
    },
    \qquad
    \gamma'
    =
    \frac{1}{\sqrt{1-\beta'^2}}.
\end{equation*}

Таким образом, при $\beta>\beta'$
%
\begin{equation*}
    \rboxed{
        \theta_{\max}
        =
        \arcsin
        \frac{
            \gamma'\beta'
        }{
            \gamma\beta
        }
    }.
\end{equation*}

Рассмотрим теперь случай
%
\begin{equation*}
    \beta<\beta'.
\end{equation*}

Можно выбрать каон, который в системе покоя $\varphi$-мезона
вылетает строго назад:
%
\begin{equation*}
    \theta'=\pi.
\end{equation*}

Для него
%
\begin{equation*}
    p'_x=-p',
    \qquad
    p'_y=0.
\end{equation*}

После преобразования Лоренца
%
\begin{equation*}
    p_x
    =
    \gamma
    \left(
        -p'
        +
        \beta\frac{E'}{c}
    \right).
\end{equation*}

Так как
%
\begin{equation*}
    p'
    =
    \beta'\frac{E'}{c},
\end{equation*}
%
имеем
%
\begin{equation*}
    p_x
    =
    \gamma\frac{E'}{c}
    (\beta-\beta').
\end{equation*}

При $\beta<\beta'$ выполняется
%
\begin{equation*}
    p_x<0,
    \qquad
    p_y=0,
\end{equation*}
%
поэтому каон в лабораторной системе может лететь строго назад:
%
\begin{equation*}
    \rboxed{
        \theta_{\max}=\pi,
        \qquad
        \beta<\beta'
    }.
\end{equation*}

При
%
\begin{equation*}
    \beta=\beta'
\end{equation*}
%
каон, вылетевший точно назад, в лабораторной системе покоится:
%
\begin{equation*}
    p_x=p_y=0.
\end{equation*}

Чтобы определить предельный угол, положим
%
\begin{equation*}
    \theta'
    =
    \pi-\delta,
    \qquad
    \delta\rightarrow0^+.
\end{equation*}

Тогда
%
\begin{equation*}
    \sin\theta'
    \simeq
    \delta,
    \qquad
    \cos\theta'
    \simeq
    -1+\frac{\delta^2}{2}.
\end{equation*}

При $\beta=\beta'$ компоненты импульса имеют вид
%
\begin{equation*}
    p_y
    \simeq
    \beta'\frac{E'}{c}\,\delta,
    \qquad
    p_x
    \simeq
    \gamma\beta'\frac{E'}{c}\,
    \frac{\delta^2}{2}.
\end{equation*}

Следовательно,
%
\begin{equation*}
    \tan\theta
    =
    \frac{p_y}{p_x}
    \simeq
    \frac{2}{\gamma\delta}
    \longrightarrow
    \infty,
\end{equation*}
%
то есть
%
\begin{equation*}
    \rboxed{
        \theta_{\max}
        \rightarrow
        \frac{\pi}{2},
        \qquad
        \beta=\beta'
    }.
\end{equation*}

Итак,
%
\begin{equation*}
    \rboxed{
        \theta_{\max}
        =
        \begin{cases}
            \pi,
            & \beta<\beta',
            \\[5pt]
            \dfrac{\pi}{2},
            & \beta=\beta',
            \\[8pt]
            \displaystyle
            \arcsin
            \frac{\gamma'\beta'}{\gamma\beta},
            & \beta>\beta'.
        \end{cases}
    }
\end{equation*}


\subsection*{Максимальный угол между каонами.}

Положим
%
\begin{equation*}
    x=\cos\theta',
    \qquad
    y=\sin\theta'
    =
    \sqrt{1-x^2},
    \qquad
    a=\frac{\beta}{\beta'}.
\end{equation*}

Для двух каонов в системе покоя $\varphi$-мезона продольные
компоненты импульса имеют противоположные знаки.
Поэтому в лабораторной системе
%
\begin{equation*}
    \tan\theta_1
    =
    \frac{
        y
    }{
        \gamma(a+x)
    },
    \qquad
    \tan\theta_2
    =
    \frac{
        y
    }{
        \gamma(a-x)
    }.
\end{equation*}

При $\beta>\beta'$ оба каона движутся в передней полусфере,
и угол между ними равен
%
\begin{equation*}
    \Psi
    =
    \theta_1+\theta_2.
\end{equation*}

Используя формулу сложения тангенсов,
%
\begin{equation*}
    \tan\Psi
    =
    \frac{
        \tan\theta_1+\tan\theta_2
    }{
        1-\tan\theta_1\tan\theta_2
    },
\end{equation*}
%
получаем
%
\begin{equation*}
    \tan\Psi
    =
    \frac{
        2a\gamma\sqrt{1-x^2}
    }{
        \gamma^2(a^2-x^2)
        -
        (1-x^2)
    }.
\end{equation*}

Для поиска максимума введём
%
\begin{equation*}
    s=x^2.
\end{equation*}

Тогда
%
\begin{equation*}
    \tan^2\Psi
    =
    \frac{
        4a^2\gamma^2(1-s)
    }{
        \left[
            \gamma^2a^2
            -
            1
            -
            (\gamma^2-1)s
        \right]^2
    }.
\end{equation*}

Условие экстремума удобно записать как
%
\begin{equation*}
    \frac{d}{ds}
    \ln\tan^2\Psi
    =
    0.
\end{equation*}

Отсюда
%
\begin{equation*}
    -\frac{1}{1-s}
    +
    \frac{
        2(\gamma^2-1)
    }{
        \gamma^2a^2
        -
        1
        -
        (\gamma^2-1)s
    }
    =
    0.
\end{equation*}

После упрощения
%
\begin{equation*}
    s_{\mathrm{ext}}
    =
    1+\frac{1}{\beta^2}-\frac{1}{\beta'^2}.
\end{equation*}

Так как $s=\cos^2\theta'$, положение максимума зависит
от величины $\beta$.

Если
%
\begin{equation*}
    s_{\mathrm{ext}}\leq0,
\end{equation*}
%
то максимум достигается на границе $s=0$, то есть
%
\begin{equation*}
    \cos\theta'=0,
    \qquad
    \theta'=\frac{\pi}{2}.
\end{equation*}

Условие $s_{\mathrm{ext}}\leq0$ можно переписать в виде
%
\begin{equation*}
    \beta
    \geq
    \gamma'\beta'.
\end{equation*}

В этом случае оба каона в системе покоя $\varphi$-мезона
вылетают поперёк направления буста. В лабораторной системе
они расположены симметрично, причём
%
\begin{equation*}
    \tan\theta_1
    =
    \tan\theta_2
    =
    \frac{
        \beta'
    }{
        \gamma\beta
    }.
\end{equation*}

Поэтому
%
\begin{equation*}
    \rboxed{
        \Psi_{\max}
        =
        2\arctan
        \frac{
            \beta'
        }{
            \gamma\beta
        },
        \qquad
        \beta\geq\gamma'\beta'
    }.
\end{equation*}

Если же
%
\begin{equation*}
    \beta'
    <
    \beta
    <
    \gamma'\beta',
\end{equation*}
%
экстремум находится внутри допустимого интервала:
%
\begin{equation*}
    \cos^2\theta'_{\mathrm{ext}}
    =
    1+\frac{1}{\beta^2}-\frac{1}{\beta'^2}.
\end{equation*}

Подставляя это значение в выражение для $\tan\Psi$, получаем
%
\begin{equation*}
    \tan\Psi_{\max}
    =
    \sqrt{
        \frac{
            1-\beta^2
        }{
            \beta^2-\beta'^2
        }
    }.
\end{equation*}

Следовательно,
%
\begin{equation*}
    \rboxed{
        \Psi_{\max}
        =
        \arctan
        \sqrt{
            \frac{
                1-\beta^2
            }{
                \beta^2-\beta'^2
            }
        },
        \qquad
        \beta'<\beta<\gamma'\beta'
    }.
\end{equation*}

При $\beta<\beta'$ можно выбрать один каон, летящий в системе
покоя $\varphi$-мезона вперёд, а второй --- строго назад.
После перехода в лабораторную систему первый остаётся направленным
вперёд, а второй --- назад. Поэтому
%
\begin{equation*}
    \rboxed{
        \Psi_{\max}
        =
        \pi,
        \qquad
        \beta<\beta'
    }.
\end{equation*}

При $\beta=\beta'$ строго задний каон покоится.
При приближении к этому случаю угол между направлениями стремится к
%
\begin{equation*}
    \rboxed{
        \Psi_{\max}
        \rightarrow
        \frac{\pi}{2},
        \qquad
        \beta=\beta'
    }.
\end{equation*}

Окончательно
%
\begin{equation*}
    \rboxed{
        \Psi_{\max}
        =
        \begin{cases}
            \pi,
            & \beta<\beta',
            \\[5pt]
            \dfrac{\pi}{2},
            & \beta=\beta',
            \\[8pt]
            \displaystyle
            \arctan
            \sqrt{
                \frac{
                    1-\beta^2
                }{
                    \beta^2-\beta'^2
                }
            },
            & \beta'<\beta<\gamma'\beta',
            \\[12pt]
            \displaystyle
            2\arctan
            \frac{
                \beta'
            }{
                \gamma\beta
            },
            & \beta\geq\gamma'\beta'.
        \end{cases}
    }
\end{equation*}

\end{document}
